\documentclass[11pt]{report}

\usepackage{amsmath}
\usepackage{amssymb}
\usepackage{amsthm}
\usepackage{mathtools}
\usepackage{multirow}

\newtheoremstyle{example} % Style name
	{10pt} % Space above
	{10pt} % Space below
	{\normalfont} % Body font
	{} % Header font
	{\bfseries} % Header font
	{.} % Separator
	{10pt} % Space after header
	{} % Header text
\theoremstyle{example}
\newtheorem{example}{Example}[section]
\newtheorem{definition}{Definition}[section]

\begin{document}

\title{FPGA Beamformer}
\author{Stewart Nash}
\date{2026}

\maketitle

\chapter{Adaptive Filtering}
\begin{center}
\textit{"Filters that adapt to a changing environment are discussed. The stochastic Wiener filtering and deterministic least-squares problems are set up, and the similarity between them is pointed out. First, a block-processing approach is taken in the solution to these two problems. Second, a fading memory (recursive) approach is taken. Then the application of these algorithms to the adaptive beamforming (ABF) problem is considered. The chapter appendix provides background material on conventional beamforming (CBF)."}
\end{center}
\vspace{1em}
\setcounter{section}{-1}
\section{Introduction}
Numerous applications exist that require a linear filtering operation. Often, the nature of that filtering task is time-varying in some nondeterministic fashion owing to nonstationarity of the underlying time series. In such situations, a filter that can adapt to a changing environment is needed. Examples include adaptive noise canceling, line enhancing, frequency tracking, and channel equalization. Beyond the discussion here, additional material on adaptive filtering can be found in references 1 through 16.

\section{The Stochastic Wiener Filtering and Deterministic Least Squares Problem}
The problem of interest is described in Figure 11.1. The goal is to filter the time series $x[n]$ with a causal, finite impulse response (FIR) digital filter to yield an estimate of the time series $s[n]$. The error in that estimate is denoted $e[n]$
\begin{equation}
	e[n]=s[n]-\hat{s}[n]=s[n]+\sum_{k=0}^N{h_kx[n-k]}
\end{equation}
Written in matrix form, Eq. 11.1 becomes
\begin{equation}
	e[n]=s[n]+[h_0\,h_1\,\dots\,h_N]
	\begin{bmatrix}
		x[n]\\
		x[n-1]\\
		\vdots\\
		x[n-N]
	\end{bmatrix}
\end{equation}
or, in vector notation,
\begin{equation}
	e[n]=s[n]+\mathbf{h}^T\mathbf{x}
\end{equation}

\end{document}
